\section{Related Work}

\subsection{Group Robustness Methods}

The worst-group accuracy problem has motivated several intervention strategies. \textbf{Group DRO} \cite{sagawa2020distributionally} reweights the loss to minimize worst-group error, achieving $\sim$91\% WGA on Waterbirds but requiring group labels during training. \textbf{Just Train Twice} \cite{liu2021just} identifies hard examples in a first training pass and upweights them in a second. \textbf{Last-layer retraining} \cite{kirichenko2023last} demonstrates that ERM learns good features---only the classifier layer needs adjustment.

Our work extends Kirichenko et al.\ \cite{kirichenko2023last} by asking \emph{why} this suppression occurs. We provide a mechanistic explanation rooted in loss landscape geometry.

\subsection{Loss Landscape Analysis}

\textbf{Sharpness-Aware Minimization (SAM)} \cite{foret2021sharpness} seeks flat minima for better generalization. Kalra \& Barkeshli \cite{kalra2023phase} characterized early training dynamics via a phase diagram identifying four regimes. Our temporal precedence finding situates the gradient-to-curvature mechanism within these early training phases.

\subsection{Feature Learning Dynamics}

\textbf{Simplicity bias} \cite{shah2020pitfalls} explains that neural networks preferentially learn simpler features first. \textbf{Shortcut learning} \cite{geirhos2020shortcut} provides a unifying framework. LaBonte \& Muthukumar \cite{labonte2026sgd} provide the first theoretical proof that SGD learns spurious features exponentially fast. Our work complements this theory with empirical characterization.
